\documentclass[
  doc,
  floatsintext,
  longtable,
  a4paper,
  nolmodern,
  notxfonts,
  notimes,
  12pt,
  colorlinks=true,linkcolor=blue,citecolor=blue,urlcolor=blue]{apa7}

\usepackage{amsmath}
\usepackage{amssymb}

\geometry{inner=1in, outer=1in}
\fancyhfoffset[LE,RO]{0cm}


\usepackage[bidi=default]{babel}
\babelprovide[main,import]{spanish}


% get rid of language-specific shorthands (see #6817):
\let\LanguageShortHands\languageshorthands
\def\languageshorthands#1{}

\RequirePackage{longtable}
\RequirePackage{threeparttablex}

\makeatletter
\renewcommand{\paragraph}{\@startsection{paragraph}{4}{\parindent}%
	{0\baselineskip \@plus 0.2ex \@minus 0.2ex}%
	{-.5em}%
	{\normalfont\normalsize\bfseries\typesectitle}}

\renewcommand{\subparagraph}[1]{\@startsection{subparagraph}{5}{0.5em}%
	{0\baselineskip \@plus 0.2ex \@minus 0.2ex}%
	{-\z@\relax}%
	{\normalfont\normalsize\bfseries\itshape\hspace{\parindent}{#1}\textit{\addperi}}{\relax}}
\makeatother




\usepackage{longtable, booktabs, multirow, multicol, colortbl, hhline, caption, array, float, xpatch}
\usepackage{subcaption}


\renewcommand\thesubfigure{\Alph{subfigure}}
\setcounter{topnumber}{2}
\setcounter{bottomnumber}{2}
\setcounter{totalnumber}{4}
\renewcommand{\topfraction}{0.85}
\renewcommand{\bottomfraction}{0.85}
\renewcommand{\textfraction}{0.15}
\renewcommand{\floatpagefraction}{0.7}

\usepackage{tcolorbox}
\tcbuselibrary{listings,theorems, breakable, skins}
\usepackage{fontawesome5}

\definecolor{quarto-callout-color}{HTML}{909090}
\definecolor{quarto-callout-note-color}{HTML}{0758E5}
\definecolor{quarto-callout-important-color}{HTML}{CC1914}
\definecolor{quarto-callout-warning-color}{HTML}{EB9113}
\definecolor{quarto-callout-tip-color}{HTML}{00A047}
\definecolor{quarto-callout-caution-color}{HTML}{FC5300}
\definecolor{quarto-callout-color-frame}{HTML}{ACACAC}
\definecolor{quarto-callout-note-color-frame}{HTML}{4582EC}
\definecolor{quarto-callout-important-color-frame}{HTML}{D9534F}
\definecolor{quarto-callout-warning-color-frame}{HTML}{F0AD4E}
\definecolor{quarto-callout-tip-color-frame}{HTML}{02B875}
\definecolor{quarto-callout-caution-color-frame}{HTML}{FD7E14}

%\newlength\Oldarrayrulewidth
%\newlength\Oldtabcolsep


\usepackage{hyperref}



\usepackage{color}
\usepackage{fancyvrb}
\newcommand{\VerbBar}{|}
\newcommand{\VERB}{\Verb[commandchars=\\\{\}]}
\DefineVerbatimEnvironment{Highlighting}{Verbatim}{commandchars=\\\{\}}
% Add ',fontsize=\small' for more characters per line
\usepackage{framed}
\definecolor{shadecolor}{RGB}{241,243,245}
\newenvironment{Shaded}{\begin{snugshade}}{\end{snugshade}}
\newcommand{\AlertTok}[1]{\textcolor[rgb]{0.68,0.00,0.00}{#1}}
\newcommand{\AnnotationTok}[1]{\textcolor[rgb]{0.37,0.37,0.37}{#1}}
\newcommand{\AttributeTok}[1]{\textcolor[rgb]{0.40,0.45,0.13}{#1}}
\newcommand{\BaseNTok}[1]{\textcolor[rgb]{0.68,0.00,0.00}{#1}}
\newcommand{\BuiltInTok}[1]{\textcolor[rgb]{0.00,0.23,0.31}{#1}}
\newcommand{\CharTok}[1]{\textcolor[rgb]{0.13,0.47,0.30}{#1}}
\newcommand{\CommentTok}[1]{\textcolor[rgb]{0.37,0.37,0.37}{#1}}
\newcommand{\CommentVarTok}[1]{\textcolor[rgb]{0.37,0.37,0.37}{\textit{#1}}}
\newcommand{\ConstantTok}[1]{\textcolor[rgb]{0.56,0.35,0.01}{#1}}
\newcommand{\ControlFlowTok}[1]{\textcolor[rgb]{0.00,0.23,0.31}{\textbf{#1}}}
\newcommand{\DataTypeTok}[1]{\textcolor[rgb]{0.68,0.00,0.00}{#1}}
\newcommand{\DecValTok}[1]{\textcolor[rgb]{0.68,0.00,0.00}{#1}}
\newcommand{\DocumentationTok}[1]{\textcolor[rgb]{0.37,0.37,0.37}{\textit{#1}}}
\newcommand{\ErrorTok}[1]{\textcolor[rgb]{0.68,0.00,0.00}{#1}}
\newcommand{\ExtensionTok}[1]{\textcolor[rgb]{0.00,0.23,0.31}{#1}}
\newcommand{\FloatTok}[1]{\textcolor[rgb]{0.68,0.00,0.00}{#1}}
\newcommand{\FunctionTok}[1]{\textcolor[rgb]{0.28,0.35,0.67}{#1}}
\newcommand{\ImportTok}[1]{\textcolor[rgb]{0.00,0.46,0.62}{#1}}
\newcommand{\InformationTok}[1]{\textcolor[rgb]{0.37,0.37,0.37}{#1}}
\newcommand{\KeywordTok}[1]{\textcolor[rgb]{0.00,0.23,0.31}{\textbf{#1}}}
\newcommand{\NormalTok}[1]{\textcolor[rgb]{0.00,0.23,0.31}{#1}}
\newcommand{\OperatorTok}[1]{\textcolor[rgb]{0.37,0.37,0.37}{#1}}
\newcommand{\OtherTok}[1]{\textcolor[rgb]{0.00,0.23,0.31}{#1}}
\newcommand{\PreprocessorTok}[1]{\textcolor[rgb]{0.68,0.00,0.00}{#1}}
\newcommand{\RegionMarkerTok}[1]{\textcolor[rgb]{0.00,0.23,0.31}{#1}}
\newcommand{\SpecialCharTok}[1]{\textcolor[rgb]{0.37,0.37,0.37}{#1}}
\newcommand{\SpecialStringTok}[1]{\textcolor[rgb]{0.13,0.47,0.30}{#1}}
\newcommand{\StringTok}[1]{\textcolor[rgb]{0.13,0.47,0.30}{#1}}
\newcommand{\VariableTok}[1]{\textcolor[rgb]{0.07,0.07,0.07}{#1}}
\newcommand{\VerbatimStringTok}[1]{\textcolor[rgb]{0.13,0.47,0.30}{#1}}
\newcommand{\WarningTok}[1]{\textcolor[rgb]{0.37,0.37,0.37}{\textit{#1}}}

\providecommand{\tightlist}{%
  \setlength{\itemsep}{0pt}\setlength{\parskip}{0pt}}
\usepackage{longtable,booktabs,array}
\newcounter{none} % for unnumbered tables
\usepackage{calc} % for calculating minipage widths
% Correct order of tables after \paragraph or \subparagraph
\usepackage{etoolbox}
\makeatletter
\patchcmd\longtable{\par}{\if@noskipsec\mbox{}\fi\par}{}{}
\makeatother
% Allow footnotes in longtable head/foot
\IfFileExists{footnotehyper.sty}{\usepackage{footnotehyper}}{\usepackage{footnote}}
\makesavenoteenv{longtable}

\usepackage{graphicx}
\makeatletter
\newsavebox\pandoc@box
\newcommand*\pandocbounded[1]{% scales image to fit in text height/width
  \sbox\pandoc@box{#1}%
  \Gscale@div\@tempa{\textheight}{\dimexpr\ht\pandoc@box+\dp\pandoc@box\relax}%
  \Gscale@div\@tempb{\linewidth}{\wd\pandoc@box}%
  \ifdim\@tempb\p@<\@tempa\p@\let\@tempa\@tempb\fi% select the smaller of both
  \ifdim\@tempa\p@<\p@\scalebox{\@tempa}{\usebox\pandoc@box}%
  \else\usebox{\pandoc@box}%
  \fi%
}
% Set default figure placement to htbp
\def\fps@figure{htbp}
\makeatother







\usepackage{newtx}

\defaultfontfeatures{Scale=MatchLowercase}
\defaultfontfeatures[\rmfamily]{Ligatures=TeX,Scale=1}





\title{Ejercicios de variables}


\shorttitle{EJERCICIOS DE VARIABLES}


\usepackage{etoolbox}



\ccoppy{\textcopyright~2026}



\author{Edison Achalma}



\affiliation{
{Escuela Profesional de Economía, Universidad Nacional de San Cristóbal
de Huamanga}}




\leftheader{Achalma}

\date{2026-10-20}


\abstract{Este abstract será actualizado una vez que se complete el
contenido final del artículo. }



\authornote{\par{\addORCIDlink{Edison Achalma}{0000-0001-6996-3364}} 
\par{ }
\par{   El autor no tiene conflictos de interés que revelar.    Los
roles de autor se clasificaron utilizando la taxonomía de roles de
colaborador (CRediT; https://credit.niso.org/) de la siguiente
manera:  Edison Achalma:   conceptualización, metodología, análisis
formal, investigación, recursos, curación de
datos, redacción, visualización, supervisión, administración del
proyecto}
\par{La correspondencia relativa a este artículo debe dirigirse a Edison
Achalma, Escuela Profesional de Economía, Universidad Nacional de San
Cristóbal de Huamanga, Portal Independencia N°
57, Ayacucho, AYA 5001, Perú, Email: \href{mailto:elmer.achalma.09@unsch.edu.pe}{elmer.achalma.09@unsch.edu.pe}}
}

\makeatletter
\let\endoldlt\endlongtable
\def\endlongtable{
\hline
\endoldlt
}
\makeatother

\urlstyle{same}



\hypersetup{
  pdftitle={Ejercicios de variables},
  pdfauthor={Autor}
}
\makeatletter
\@ifpackageloaded{caption}{}{\usepackage{caption}}
\AtBeginDocument{%
\ifdefined\contentsname
  \renewcommand*\contentsname{Tabla de contenidos}
\else
  \newcommand\contentsname{Tabla de contenidos}
\fi
\ifdefined\listfigurename
  \renewcommand*\listfigurename{Lista de Figuras}
\else
  \newcommand\listfigurename{Lista de Figuras}
\fi
\ifdefined\listtablename
  \renewcommand*\listtablename{Lista de Tablas}
\else
  \newcommand\listtablename{Lista de Tablas}
\fi
\ifdefined\figurename
  \renewcommand*\figurename{Figura}
\else
  \newcommand\figurename{Figura}
\fi
\ifdefined\tablename
  \renewcommand*\tablename{Tabla}
\else
  \newcommand\tablename{Tabla}
\fi
}
\@ifpackageloaded{float}{}{\usepackage{float}}
\floatstyle{ruled}
\@ifundefined{c@chapter}{\newfloat{codelisting}{h}{lop}}{\newfloat{codelisting}{h}{lop}[chapter]}
\floatname{codelisting}{Listado}
\newcommand*\listoflistings{\listof{codelisting}{Lista de Listados}}
\makeatother
\makeatletter
\makeatother
\makeatletter
\@ifpackageloaded{caption}{}{\usepackage{caption}}
\@ifpackageloaded{subcaption}{}{\usepackage{subcaption}}
\makeatother
\makeatletter
\@ifpackageloaded{fontawesome5}{}{\usepackage{fontawesome5}}
\makeatother

% From https://tex.stackexchange.com/a/645996/211326
%%% apa7 doesn't want to add appendix section titles in the toc
%%% let's make it do it
\makeatletter
\xpatchcmd{\appendix}
  {\par}
  {\addcontentsline{toc}{section}{\@currentlabelname}\par}
  {}{}
\makeatother

%% Disable longtable counter
%% https://tex.stackexchange.com/a/248395/211326

\usepackage{etoolbox}

\makeatletter
\patchcmd{\LT@caption}
  {\bgroup}
  {\bgroup\global\LTpatch@captiontrue}
  {}{}
\patchcmd{\longtable}
  {\par}
  {\par\global\LTpatch@captionfalse}
  {}{}
\apptocmd{\endlongtable}
  {\ifLTpatch@caption\else\addtocounter{table}{-1}\fi}
  {}{}
\newif\ifLTpatch@caption
\makeatother

\begin{document}

\maketitle


\hypertarget{toc}{}
\tableofcontents
\newpage
\section[Introduction]{Ejercicios de variables}

\setcounter{secnumdepth}{3}

\setlength\LTleft{0pt}




¡Excelente! Has completado la teoría sobre variables en Python. Ahora es
momento de poner a prueba tus conocimientos con estos ejercicios
prácticos. Te recomiendo resolverlos en orden, ya que aumentan
gradualmente en dificultad.

\begin{center}\rule{0.5\linewidth}{0.5pt}\end{center}

\subsection{📌 Categoría 1: Variables
Básicas}\label{categoruxeda-1-variables-buxe1sicas}

\subsubsection{Ejercicio 1}\label{ejercicio-1}

\textbf{Completa el código para crear una variable llamada
\texttt{ciudad} con el valor ``Madrid'':}

\begin{Shaded}
\begin{Highlighting}[]
\CommentTok{\# Escribe tu código aquí}
\end{Highlighting}
\end{Shaded}

💡 Ver solución

\begin{Shaded}
\begin{Highlighting}[]
\NormalTok{ciudad }\OperatorTok{=} \StringTok{"Madrid"}
\end{Highlighting}
\end{Shaded}

\begin{center}\rule{0.5\linewidth}{0.5pt}\end{center}

\subsubsection{Ejercicio 2}\label{ejercicio-2}

\textbf{¿Cuál es la salida del siguiente código?}

\begin{Shaded}
\begin{Highlighting}[]
\NormalTok{x }\OperatorTok{=} \DecValTok{10}
\NormalTok{y }\OperatorTok{=} \DecValTok{20}
\NormalTok{x }\OperatorTok{=}\NormalTok{ y}
\BuiltInTok{print}\NormalTok{(x)}
\end{Highlighting}
\end{Shaded}

💡 Ver solución

\textbf{Salida:} \texttt{20}

\textbf{Explicación:} \texttt{x\ =\ y} asigna el valor de \texttt{y}
(que es 20) a \texttt{x}, sobrescribiendo el valor anterior de 10.

\begin{center}\rule{0.5\linewidth}{0.5pt}\end{center}

\subsubsection{Ejercicio 3}\label{ejercicio-3}

\textbf{Corrige el error en el siguiente código:}

\begin{Shaded}
\begin{Highlighting}[]
\DecValTok{2}\ErrorTok{nombre} \OperatorTok{=} \StringTok{"Ana"}
\BuiltInTok{print}\NormalTok{(}\DecValTok{2}\ErrorTok{nombre}\NormalTok{)}
\end{Highlighting}
\end{Shaded}

💡 Ver solución

\begin{Shaded}
\begin{Highlighting}[]
\NormalTok{nombre2 }\OperatorTok{=} \StringTok{"Ana"}  \CommentTok{\# Las variables no pueden empezar con número}
\BuiltInTok{print}\NormalTok{(nombre2)}
\end{Highlighting}
\end{Shaded}

\begin{center}\rule{0.5\linewidth}{0.5pt}\end{center}

\subsection{📌 Categoría 2: Nombres de
Variables}\label{categoruxeda-2-nombres-de-variables}

\subsubsection{Ejercicio 4}\label{ejercicio-4}

\textbf{¿Cuáles de estos nombres de variables son válidos en Python?}
(Selecciona todos los correctos)

\begin{itemize}
\tightlist
\item
  \begin{enumerate}
  \def\labelenumi{\Alph{enumi})}
  \tightlist
  \item
    \texttt{mi-variable}
  \end{enumerate}
\item
  \begin{enumerate}
  \def\labelenumi{\Alph{enumi})}
  \setcounter{enumi}{1}
  \tightlist
  \item
    \texttt{\_variable}
  \end{enumerate}
\item
  \begin{enumerate}
  \def\labelenumi{\Alph{enumi})}
  \setcounter{enumi}{2}
  \tightlist
  \item
    \texttt{1variable}
  \end{enumerate}
\item
  \begin{enumerate}
  \def\labelenumi{\Alph{enumi})}
  \setcounter{enumi}{3}
  \tightlist
  \item
    \texttt{variable\_1}
  \end{enumerate}
\item
  \begin{enumerate}
  \def\labelenumi{\Alph{enumi})}
  \setcounter{enumi}{4}
  \tightlist
  \item
    \texttt{mi\ variable}
  \end{enumerate}
\end{itemize}

💡 Ver solución

\textbf{Respuestas correctas: B y D}

\begin{itemize}
\item
  \texttt{\_variable} ✅ (empieza con guión bajo, válido)
\item
  \texttt{variable\_1} ✅ (contiene letras, números y guión bajo)
\item
  \texttt{mi-variable} ❌ (contiene guión \texttt{-})
\item
  \texttt{1variable} ❌ (empieza con número)
\item
  \texttt{mi\ variable} ❌ (contiene espacio)
\end{itemize}

\begin{center}\rule{0.5\linewidth}{0.5pt}\end{center}

\subsubsection{Ejercicio 5}\label{ejercicio-5}

\textbf{Convierte estos nombres a la convención snake\_case (recomendada
en Python):}

\begin{Shaded}
\begin{Highlighting}[]
\CommentTok{\# ❌ Malos nombres}
\NormalTok{nombreCompleto }\OperatorTok{=} \StringTok{"Ana"}
\NormalTok{miVar }\OperatorTok{=} \DecValTok{25}
\NormalTok{precioTotalProducto }\OperatorTok{=} \FloatTok{99.99}
\end{Highlighting}
\end{Shaded}

💡 Ver solución

\begin{Shaded}
\begin{Highlighting}[]
\CommentTok{\# ✅ Snake case (recomendado)}
\NormalTok{nombre\_completo }\OperatorTok{=} \StringTok{"Ana"}
\NormalTok{mi\_var }\OperatorTok{=} \DecValTok{25}
\NormalTok{precio\_total\_producto }\OperatorTok{=} \FloatTok{99.99}
\end{Highlighting}
\end{Shaded}

\begin{center}\rule{0.5\linewidth}{0.5pt}\end{center}

\subsection{📌 Categoría 3: Múltiples
Valores}\label{categoruxeda-3-muxfaltiples-valores}

\subsubsection{Ejercicio 6}\label{ejercicio-6}

\textbf{Completa el código para asignar los valores correctos a las
variables:}

\begin{Shaded}
\begin{Highlighting}[]
\CommentTok{\# Asigna "Perro", "Gato", "Pájaro" a las variables a, b, c en una línea}
\CommentTok{\# Escribe tu código aquí}

\BuiltInTok{print}\NormalTok{(a)  }\CommentTok{\# Debe imprimir: Perro}
\BuiltInTok{print}\NormalTok{(b)  }\CommentTok{\# Debe imprimir: Gato}
\BuiltInTok{print}\NormalTok{(c)  }\CommentTok{\# Debe imprimir: Pájaro}
\end{Highlighting}
\end{Shaded}

💡 Ver solución

\begin{Shaded}
\begin{Highlighting}[]
\NormalTok{a, b, c }\OperatorTok{=} \StringTok{"Perro"}\NormalTok{, }\StringTok{"Gato"}\NormalTok{, }\StringTok{"Pájaro"}
\end{Highlighting}
\end{Shaded}

\begin{center}\rule{0.5\linewidth}{0.5pt}\end{center}

\subsubsection{Ejercicio 7}\label{ejercicio-7}

\textbf{¿Qué imprime el siguiente código?}

\begin{Shaded}
\begin{Highlighting}[]
\NormalTok{colores }\OperatorTok{=}\NormalTok{ [}\StringTok{"rojo"}\NormalTok{, }\StringTok{"verde"}\NormalTok{, }\StringTok{"azul"}\NormalTok{]}
\NormalTok{x, y, z }\OperatorTok{=}\NormalTok{ colores}
\BuiltInTok{print}\NormalTok{(y)}
\end{Highlighting}
\end{Shaded}

💡 Ver solución

\textbf{Salida:} \texttt{verde}

\textbf{Explicación:} El desempaquetado asigna \texttt{x\ =\ "rojo"},
\texttt{y\ =\ "verde"}, \texttt{z\ =\ "azul"}.

\begin{center}\rule{0.5\linewidth}{0.5pt}\end{center}

\subsubsection{Ejercicio 8}\label{ejercicio-8}

\textbf{Completa el código para intercambiar los valores de \texttt{x} e
\texttt{y} en una línea:}

\begin{Shaded}
\begin{Highlighting}[]
\NormalTok{x }\OperatorTok{=} \DecValTok{5}
\NormalTok{y }\OperatorTok{=} \DecValTok{10}

\CommentTok{\# Escribe tu código aquí (1 línea)}

\BuiltInTok{print}\NormalTok{(x)  }\CommentTok{\# Debe imprimir: 10}
\BuiltInTok{print}\NormalTok{(y)  }\CommentTok{\# Debe imprimir: 5}
\end{Highlighting}
\end{Shaded}

💡 Ver solución

\begin{Shaded}
\begin{Highlighting}[]
\NormalTok{x, y }\OperatorTok{=}\NormalTok{ y, x  }\CommentTok{\# Intercambio en una línea}
\end{Highlighting}
\end{Shaded}

\begin{center}\rule{0.5\linewidth}{0.5pt}\end{center}

\subsection{📌 Categoría 4: Salida de
Variables}\label{categoruxeda-4-salida-de-variables}

\subsubsection{Ejercicio 9}\label{ejercicio-9}

\textbf{¿Cuál es la forma correcta de imprimir el texto ``Ana tiene 25
años''?}

\begin{Shaded}
\begin{Highlighting}[]
\NormalTok{nombre }\OperatorTok{=} \StringTok{"Ana"}
\NormalTok{edad }\OperatorTok{=} \DecValTok{25}
\end{Highlighting}
\end{Shaded}

\textbf{Opciones:}

\begin{itemize}
\tightlist
\item
  \begin{enumerate}
  \def\labelenumi{\Alph{enumi})}
  \tightlist
  \item
    \texttt{print(nombre\ +\ "\ tiene\ "\ +\ edad\ +\ "\ años")}
  \end{enumerate}
\item
  \begin{enumerate}
  \def\labelenumi{\Alph{enumi})}
  \setcounter{enumi}{1}
  \tightlist
  \item
    \texttt{print(nombre,\ "tiene",\ edad,\ "años")}
  \end{enumerate}
\item
  \begin{enumerate}
  \def\labelenumi{\Alph{enumi})}
  \setcounter{enumi}{2}
  \tightlist
  \item
    Ambas son correctas
  \end{enumerate}
\end{itemize}

💡 Ver solución

\textbf{Respuesta correcta: B}

\textbf{Explicación:}

\begin{itemize}
\item
  \begin{enumerate}
  \def\labelenumi{\Alph{enumi})}
  \tightlist
  \item
    ❌ Error: no puedes sumar texto (\texttt{nombre}) con número
    (\texttt{edad}) usando \texttt{+}
  \end{enumerate}
\item
  \begin{enumerate}
  \def\labelenumi{\Alph{enumi})}
  \setcounter{enumi}{1}
  \tightlist
  \item
    ✅ Correcto: las comas separan diferentes tipos sin problemas
  \end{enumerate}
\end{itemize}

\begin{center}\rule{0.5\linewidth}{0.5pt}\end{center}

\subsubsection{Ejercicio 10}\label{ejercicio-10}

\textbf{¿Qué imprime el siguiente código?}

\begin{Shaded}
\begin{Highlighting}[]
\NormalTok{x }\OperatorTok{=} \DecValTok{10}
\NormalTok{y }\OperatorTok{=} \StringTok{"20"}
\BuiltInTok{print}\NormalTok{(x }\OperatorTok{+} \BuiltInTok{int}\NormalTok{(y))}
\end{Highlighting}
\end{Shaded}

💡 Ver solución

\textbf{Salida:} \texttt{30}

\textbf{Explicación:} \texttt{int(y)} convierte el texto ``20'' al
número 20, y luego suma 10 + 20 = 30.

\begin{center}\rule{0.5\linewidth}{0.5pt}\end{center}

\subsubsection{Ejercicio 11}\label{ejercicio-11}

\textbf{Completa el código para imprimir los datos separados por ''
\textbar{} ``:}

\begin{Shaded}
\begin{Highlighting}[]
\NormalTok{nombre }\OperatorTok{=} \StringTok{"Ana"}
\NormalTok{edad }\OperatorTok{=} \DecValTok{25}
\NormalTok{ciudad }\OperatorTok{=} \StringTok{"Madrid"}

\CommentTok{\# Escribe tu código aquí (usando sep)}
\CommentTok{\# Debe imprimir: Ana | 25 | Madrid}
\end{Highlighting}
\end{Shaded}

💡 Ver solución

\begin{Shaded}
\begin{Highlighting}[]
\BuiltInTok{print}\NormalTok{(nombre, edad, ciudad, sep}\OperatorTok{=}\StringTok{" | "}\NormalTok{)}
\end{Highlighting}
\end{Shaded}

\begin{center}\rule{0.5\linewidth}{0.5pt}\end{center}

\subsection{📌 Categoría 5: Variables
Globales}\label{categoruxeda-5-variables-globales}

\subsubsection{Ejercicio 12}\label{ejercicio-12}

\textbf{¿Qué imprime el siguiente código?}

\begin{Shaded}
\begin{Highlighting}[]
\NormalTok{x }\OperatorTok{=} \StringTok{"global"}

\KeywordTok{def}\NormalTok{ funcion():}
\NormalTok{    x }\OperatorTok{=} \StringTok{"local"}
    \BuiltInTok{print}\NormalTok{(x)}

\NormalTok{funcion()}
\BuiltInTok{print}\NormalTok{(x)}
\end{Highlighting}
\end{Shaded}

💡 Ver solución

\textbf{Salida:}

\begin{verbatim}
local
global
\end{verbatim}

\textbf{Explicación:} Dentro de \texttt{funcion()},
\texttt{x\ =\ "local"} crea una variable local que solo existe dentro de
la función. La variable global \texttt{x\ =\ "global"} permanece sin
cambios.

\begin{center}\rule{0.5\linewidth}{0.5pt}\end{center}

\subsubsection{Ejercicio 13}\label{ejercicio-13}

\textbf{Completa el código para modificar la variable global
\texttt{contador} desde la función:}

\begin{Shaded}
\begin{Highlighting}[]
\NormalTok{contador }\OperatorTok{=} \DecValTok{0}

\KeywordTok{def}\NormalTok{ incrementar():}
    \CommentTok{\# Escribe tu código aquí}
\NormalTok{    contador }\OperatorTok{+=} \DecValTok{1}

\NormalTok{incrementar()}
\BuiltInTok{print}\NormalTok{(contador)  }\CommentTok{\# Debe imprimir: 1}
\end{Highlighting}
\end{Shaded}

💡 Ver solución

\begin{Shaded}
\begin{Highlighting}[]
\NormalTok{contador }\OperatorTok{=} \DecValTok{0}

\KeywordTok{def}\NormalTok{ incrementar():}
    \KeywordTok{global}\NormalTok{ contador  }\CommentTok{\# Declara que vas a modificar la global}
\NormalTok{    contador }\OperatorTok{+=} \DecValTok{1}

\NormalTok{incrementar()}
\BuiltInTok{print}\NormalTok{(contador)  }\CommentTok{\# 1}
\end{Highlighting}
\end{Shaded}

\begin{center}\rule{0.5\linewidth}{0.5pt}\end{center}

\subsubsection{Ejercicio 14}\label{ejercicio-14}

\textbf{¿Qué imprime el siguiente código?}

\begin{Shaded}
\begin{Highlighting}[]
\NormalTok{mensaje }\OperatorTok{=} \StringTok{"Hola"}

\KeywordTok{def}\NormalTok{ cambiar\_mensaje():}
    \KeywordTok{global}\NormalTok{ mensaje}
\NormalTok{    mensaje }\OperatorTok{=} \StringTok{"Adiós"}
    \BuiltInTok{print}\NormalTok{(mensaje)}

\NormalTok{cambiar\_mensaje()}
\BuiltInTok{print}\NormalTok{(mensaje)}
\end{Highlighting}
\end{Shaded}

💡 Ver solución

\textbf{Salida:}

\begin{verbatim}
Adiós
Adiós
\end{verbatim}

\textbf{Explicación:} El \texttt{global} dentro de la función permite
modificar la variable global \texttt{mensaje}. Después de la función, la
global sigue siendo ``Adiós''.

\begin{center}\rule{0.5\linewidth}{0.5pt}\end{center}

\subsection{🏆 Ejercicio Final}\label{ejercicio-final}

\subsubsection{Ejercicio 15 (Integrador)}\label{ejercicio-15-integrador}

\textbf{Escribe un programa completo que:}

\begin{enumerate}
\def\labelenumi{\arabic{enumi}.}
\tightlist
\item
  Cree una variable global \texttt{total\_ventas\ =\ 0}
\item
  Cree una función \texttt{registrar\_venta(monto)} que:

  \begin{itemize}
  \tightlist
  \item
    Use la variable global para sumar el monto
  \item
    Muestre ``Venta registrada: {[}monto{]} €''
  \end{itemize}
\item
  Cree una función \texttt{mostrar\_total()} que muestre ``Total de
  ventas: {[}total{]} €''
\item
  Registre tres ventas de diferentes montos
\item
  Muestre el total final
\end{enumerate}

\textbf{Usa nombres de variables descriptivos en snake\_case.}

\begin{center}\rule{0.5\linewidth}{0.5pt}\end{center}

💡 Ver solución

\begin{Shaded}
\begin{Highlighting}[]
\CommentTok{\# Variable global}
\NormalTok{total\_ventas }\OperatorTok{=} \DecValTok{0}

\KeywordTok{def}\NormalTok{ registrar\_venta(monto):}
    \CommentTok{"""Registra una venta y actualiza el total."""}
    \KeywordTok{global}\NormalTok{ total\_ventas}
\NormalTok{    total\_ventas }\OperatorTok{+=}\NormalTok{ monto}
    \BuiltInTok{print}\NormalTok{(}\SpecialStringTok{f"Venta registrada: }\SpecialCharTok{\{}\NormalTok{monto}\SpecialCharTok{\}}\SpecialStringTok{ €"}\NormalTok{)}

\KeywordTok{def}\NormalTok{ mostrar\_total():}
    \CommentTok{"""Muestra el total de ventas acumulado."""}
    \BuiltInTok{print}\NormalTok{(}\SpecialStringTok{f"Total de ventas: }\SpecialCharTok{\{}\NormalTok{total\_ventas}\SpecialCharTok{\}}\SpecialStringTok{ €"}\NormalTok{)}

\CommentTok{\# Registrar ventas}
\NormalTok{registrar\_venta(}\FloatTok{150.50}\NormalTok{)}
\NormalTok{registrar\_venta(}\FloatTok{75.25}\NormalTok{)}
\NormalTok{registrar\_venta(}\FloatTok{200.00}\NormalTok{)}

\CommentTok{\# Mostrar resultado final}
\NormalTok{mostrar\_total()}
\end{Highlighting}
\end{Shaded}

\textbf{Salida esperada:}

\begin{verbatim}
Venta registrada: 150.5 €
Venta registrada: 75.25 €
Venta registrada: 200.0 €
Total de ventas: 425.75 €
\end{verbatim}

\begin{center}\rule{0.5\linewidth}{0.5pt}\end{center}

\subsection{📊 Resumen de Tu Progreso}\label{resumen-de-tu-progreso}

{\def\LTcaptype{none} % do not increment counter
\begin{longtable}[]{@{}lll@{}}
\toprule\noalign{}
Categoría & Ejercicios & Completados \\
\midrule\noalign{}
\endhead
\bottomrule\noalign{}
\endlastfoot
Variables Básicas & 1-3 & {[} {]} {[} {]} {[} {]} \\
Nombres de Variables & 4-5 & {[} {]} {[} {]} \\
Múltiples Valores & 6-8 & {[} {]} {[} {]} {[} {]} \\
Salida de Variables & 9-11 & {[} {]} {[} {]} {[} {]} \\
Variables Globales & 12-14 & {[} {]} {[} {]} {[} {]} \\
Desafío Final & 15 & {[} {]} \\
\end{longtable}
}

\section{Publicaciones Similares}\label{publicaciones-similares}

Si te interesó este artículo, te recomendamos que explores otros blogs y
recursos relacionados que pueden ampliar tus conocimientos. Aquí te dejo
algunas sugerencias:

\begin{enumerate}
\def\labelenumi{\arabic{enumi}.}
\tightlist
\item
  \href{https://numerus-scriptum.netlify.app/python/2020-06-19-instalacion-de-anaconda/index.pdf}{\faIcon{file-pdf}}
  \href{https://numerus-scriptum.netlify.app/python/2020-06-19-instalacion-de-anaconda}{Instalacion
  De Anaconda}
\item
  \href{https://numerus-scriptum.netlify.app/python/2020-06-20-configurar-entorno-virtual-python-anaconda/index.pdf}{\faIcon{file-pdf}}
  \href{https://numerus-scriptum.netlify.app/python/2020-06-20-configurar-entorno-virtual-python-anaconda}{Configurar
  Entorno Virtual Python Anaconda}
\item
  \href{https://numerus-scriptum.netlify.app/python/2021-04-17-01-introducion-a-la-programacion-con-python/index.pdf}{\faIcon{file-pdf}}
  \href{https://numerus-scriptum.netlify.app/python/2021-04-17-01-introducion-a-la-programacion-con-python}{01
  Introducion A La Programacion Con Python}
\item
  \href{https://numerus-scriptum.netlify.app/python/2021-05-31-02-variables-expresiones-y-statements-con-python/index.pdf}{\faIcon{file-pdf}}
  \href{https://numerus-scriptum.netlify.app/python/2021-05-31-02-variables-expresiones-y-statements-con-python}{02
  Variables Expresiones Y Statements Con Python}
\item
  \href{https://numerus-scriptum.netlify.app/python/2021-06-07-03-objetos-de-python/index.pdf}{\faIcon{file-pdf}}
  \href{https://numerus-scriptum.netlify.app/python/2021-06-07-03-objetos-de-python}{03
  Objetos De Python}
\item
  \href{https://numerus-scriptum.netlify.app/python/2021-06-14-04-ejecucion-condicional-con-python/index.pdf}{\faIcon{file-pdf}}
  \href{https://numerus-scriptum.netlify.app/python/2021-06-14-04-ejecucion-condicional-con-python}{04
  Ejecucion Condicional Con Python}
\item
  \href{https://numerus-scriptum.netlify.app/python/2021-06-21-05-iteraciones-con-python/index.pdf}{\faIcon{file-pdf}}
  \href{https://numerus-scriptum.netlify.app/python/2021-06-21-05-iteraciones-con-python}{05
  Iteraciones Con Python}
\item
  \href{https://numerus-scriptum.netlify.app/python/2021-08-16-06-funciones-con-python/index.pdf}{\faIcon{file-pdf}}
  \href{https://numerus-scriptum.netlify.app/python/2021-08-16-06-funciones-con-python}{06
  Funciones Con Python}
\item
  \href{https://numerus-scriptum.netlify.app/python/2021-08-23-07-dataframes-con-python/index.pdf}{\faIcon{file-pdf}}
  \href{https://numerus-scriptum.netlify.app/python/2021-08-23-07-dataframes-con-python}{07
  Dataframes Con Python}
\item
  \href{https://numerus-scriptum.netlify.app/python/2021-11-29-08-prediccion-y-metrica-de-performance-con-python/index.pdf}{\faIcon{file-pdf}}
  \href{https://numerus-scriptum.netlify.app/python/2021-11-29-08-prediccion-y-metrica-de-performance-con-python}{08
  Prediccion Y Metrica De Performance Con Python}
\item
  \href{https://numerus-scriptum.netlify.app/python/2021-12-06-09-metodos-de-machine-learning-para-clasificacion-con-python/index.pdf}{\faIcon{file-pdf}}
  \href{https://numerus-scriptum.netlify.app/python/2021-12-06-09-metodos-de-machine-learning-para-clasificacion-con-python}{09
  Metodos De Machine Learning Para Clasificacion Con Python}
\item
  \href{https://numerus-scriptum.netlify.app/python/2021-12-13-10-metodos-de-machine-learning-para-regresion-con-python/index.pdf}{\faIcon{file-pdf}}
  \href{https://numerus-scriptum.netlify.app/python/2021-12-13-10-metodos-de-machine-learning-para-regresion-con-python}{10
  Metodos De Machine Learning Para Regresion Con Python}
\item
  \href{https://numerus-scriptum.netlify.app/python/2025-05-10-visualizacion-de-datos-con-python/index.pdf}{\faIcon{file-pdf}}
  \href{https://numerus-scriptum.netlify.app/python/2025-05-10-visualizacion-de-datos-con-python}{Visualizacion
  De Datos Con Python}
\end{enumerate}

Esperamos que encuentres estas publicaciones igualmente interesantes y
útiles. ¡Disfruta de la lectura!






\end{document}
